\section{Exercícios N1--N4}
\begin{frame}{N1 — fundamentos industriais}
\small
\begin{enumerate}
\item Diferencie QGBT, QD, CCM e painel de automação.
\item Explique as funções de contator, relé de sobrecarga, disjuntor-motor, soft-starter e VFD.
\item Calcule a corrente nominal de um motor trifásico de \SI{15}{kW}, \SI{380}{V}, $fp=0{,}86$ e $\eta=0{,}90$.
\item Descreva a função da equipotencialização em uma planta industrial.
\end{enumerate}
\end{frame}

\begin{frame}{N2 — dimensionamento industrial}
\small
\begin{enumerate}
\item Dimensione preliminarmente o alimentador de um motor de \SI{30}{kW} a \SI{55}{m} do CCM.
\item Compare partida direta, estrela-triângulo, soft-starter e VFD para uma carga de alto torque de partida.
\item Verifique queda de tensão de um alimentador trifásico e proponha duas correções possíveis.
\item Defina critérios para seleção da capacidade de interrupção de um disjuntor industrial.
\end{enumerate}
\end{frame}

\begin{frame}{N3 — análise industrial}
\small
\begin{enumerate}
\item Estime a corrente de curto-circuito no secundário de um transformador a partir de sua potência e impedância percentual.
\item Discuta como obter seletividade entre proteção do QGBT e do CCM.
\item Avalie os efeitos de VFDs sobre harmônicas, fator de potência e corrente no neutro.
\item Proponha ações para reduzir risco de arco elétrico durante manutenção.
\end{enumerate}
\end{frame}

\begin{frame}{N4 — desafio de projeto industrial}
\small
Uma planta possui transformador próprio, QGBT, dois CCMs, motores de \SI{45}{kW}, \SI{30}{kW} e \SI{18.5}{kW}, cargas de soldagem e painéis de automação.
\begin{enumerate}
\item Estruture a arquitetura elétrica e o diagrama unifilar.
\item Defina critérios de demanda, curto-circuito e seletividade.
\item Proponha alimentação, proteção e acionamento dos motores.
\item Inclua aterramento, qualidade de energia, NR-10 e plano de manutenção.
\end{enumerate}
\end{frame}

